%==============================================================================
\section{Toward Iterative Optimization}
\label{sec:optimization}
%==============================================================================

A key design goal of the evaluation framework is to produce data that directly enables prompt optimization. In this section, we describe the optimization methodology that the framework supports and present preliminary results from an initial optimization experiment.

\subsection{Optimization Loop}

The evaluation framework enables a closed-loop optimization cycle:

\begin{enumerate}
    \item \textbf{Generate.} Run \threatforest{} across evaluation domains, producing traced outputs.
    \item \textbf{Score.} SMEs review outputs using the framework's scoring dimensions; automated metrics are computed in parallel.
    \item \textbf{Export.} The export pipeline produces Langfuse Dataset items with input/expected\_output pairs.
    \item \textbf{Optimize.} DSPy's MIPROv2~\cite{opsahlong2024miprov2} jointly optimizes prompt instructions and few-shot demonstrations using Bayesian search over the exported dataset.
    \item \textbf{Evaluate.} The optimized prompt is evaluated against a held-out test set using the same framework metrics.
    \item \textbf{Deploy.} If improved, the optimized prompt replaces the baseline; the cycle repeats.
\end{enumerate}

\subsection{MIPROv2 Integration}

We configure MIPROv2 with three optimization modes trading off thoroughness for compute cost:

\begin{table}[t]
\centering
\caption{MIPROv2 optimization configurations.}
\label{tab:mipro}
\begin{tabular}{@{}lrrr@{}}
\toprule
\textbf{Mode} & \textbf{Trials} & \textbf{Candidates} & \textbf{Est.\ Duration} \\
\midrule
Light & 5--10 & 5 & 5--10 min \\
Medium & 15--20 & 10 & 30--60 min \\
Heavy & 30+ & 15 & 1--2 hr \\
\bottomrule
\end{tabular}
\end{table}

The optimizer receives the evaluation framework's composite score as its objective function, with the weighted metric decomposition (syntax 30\%, path coverage 30\%, color coding 20\%, logical flow 20\%) providing gradient signal across quality dimensions.

\subsection{Preliminary Experiment}

We conducted an initial optimization experiment on attack tree generation using a synthetic dataset of 7 examples with a 5/2 train/test split:

\begin{table}[t]
\centering
\caption{Preliminary optimization results (attack tree generation, $n=7$).}
\label{tab:prelim}
\begin{tabular}{@{}lcc@{}}
\toprule
\textbf{Metric} & \textbf{Baseline} & \textbf{Optimized} \\
\midrule
Average Score & 0.500 & 0.500 \\
Success Rate & 50.0\% & 50.0\% \\
\bottomrule
\end{tabular}
\end{table}

The optimizer did not improve over baseline, which we attribute to the extremely small dataset size (7 examples, well below the 20+ recommended by DSPy documentation). This result underscores the importance of the evaluation framework's data collection infrastructure: meaningful optimization requires a critical mass of SME-scored examples that can only be accumulated through systematic evaluation.

\subsection{Roadmap}

With the evaluation framework in place, the path to optimization is:

\begin{enumerate}
    \item \textbf{Baseline collection (this work).} Run \threatforest{} across all evaluation domains and collect SME scores using the framework.
    \item \textbf{Dataset expansion.} Accumulate 50--100 SME-scored examples across domains via the tracing infrastructure.
    \item \textbf{Evaluator calibration.} Train an automated evaluator (LLM-as-Judge) against SME scores, targeting correlation $r > 0.85$.
    \item \textbf{Full optimization.} Run MIPROv2 in heavy mode with the expanded dataset.
    \item \textbf{Multi-stage optimization.} Extend optimization beyond attack trees to threat statement generation and TTP matching prompts.
\end{enumerate}
